\documentclass{article}

\usepackage{kotex}
\usepackage{amsmath}
\usepackage{amssymb}
\usepackage{booktabs}
\usepackage{geometry}
\geometry{a4paper, margin=2.5cm}

\begin{document}

\subsection*{COLREGs Compliance Metric}

COLREGs 준수도는 조우(encounter) 단위로 평가함.
두 선박이 COLREGs 상황(Head-on, Crossing, Overtaking)에 연속 5스텝 이상 놓인 구간을 하나의 조우 $k$로 정의하고, 해당 시간 구간을 $\mathcal{T}_k$로 표기함.

조우 구간에서 정규화된 타각 $\bar{\delta}_t = \delta_t / \delta_{\max} \in [-1,\;1]$을 매 스텝 기록하여 다음 세 지표를 산출함.

\begin{equation}
  R_{\text{stb}} = \frac{1}{|\mathcal{T}_k|}\sum_{t \in \mathcal{T}_k} \mathbf{1}\!\left[\bar{\delta}_t > \epsilon_s\right], \qquad
  R_{\text{port}} = \frac{1}{|\mathcal{T}_k|}\sum_{t \in \mathcal{T}_k} \mathbf{1}\!\left[\bar{\delta}_t < -\epsilon_p\right]
\end{equation}

\begin{equation}
  R_{\text{ck}} = \frac{1}{|\mathcal{T}_k|}\sum_{t \in \mathcal{T}_k} \mathbf{1}\!\left[|\bar{\delta}_t| < \epsilon_c\right]
\end{equation}

\noindent $R_{\text{stb}}$는 우현 변침 비율, $R_{\text{port}}$는 좌현 위반 비율, $R_{\text{ck}}$는 침로 유지 비율이며, 임계값은 $\epsilon_s{=}0.05$, $\epsilon_p{=}0.1$, $\epsilon_c{=}0.15$로 설정함.

각 조우의 준수 여부는 상황 유형에 따라 Table~\ref{tab:colregs_criteria}의 조건으로 판정함.

\begin{table}[h]
\centering
\caption{COLREGs 상황별 준수 판정 기준. GW: Give-Way, SO: Stand-On.}
\label{tab:colregs_criteria}
\renewcommand{\arraystretch}{1.3}
\begin{tabular}{@{}lcl@{}}
\toprule
\textbf{상황} & \textbf{규칙} & \textbf{준수 조건} \\
\midrule
Head-on       & Rule 14    & $R_{\text{stb}} > 0.5 \;\wedge\; R_{\text{port}} < 0.2$ \\
Crossing (GW) & Rule 15/16 & $R_{\text{stb}} > 0.5 \;\wedge\; R_{\text{port}} < 0.2$ \\
Crossing (SO) & Rule 17    & $R_{\text{ck}} > 0.6$ \\
Overtaking    & Rule 13    & $R_{\text{stb}} > 0.3$ \\
\bottomrule
\end{tabular}
\end{table}

\noindent 전체 준수율은 준수 조우 수를 전체 조우 수로 나눈 비율로 산출하며, 조우 중 선박 쌍 간 최소 통과 거리를 DCPA 근사값으로 함께 기록함.

\begin{equation}
  C = \frac{N_{\text{compliant}}}{N_{\text{total}}}
\end{equation}

\end{document}
